\documentclass[10.5pt]{article}
\usepackage[margin=0.85in, top=0.75in, bottom=0.75in]{geometry}
\usepackage{microtype}
\usepackage{parskip}
\usepackage[colorlinks=true,linkcolor=black,citecolor=black,urlcolor=blue]{hyperref}
\usepackage{booktabs}
\usepackage{listings}
\lstset{
  basicstyle=\ttfamily\footnotesize,
  breaklines=true,
  frame=single,
  framesep=4pt,
  xleftmargin=6pt,
  xrightmargin=6pt
}

\title{\textbf{Using an LLM in a Formal Verification Loop to Discover\\[4pt]
Detailed-Balance-Satisfying Monte Carlo Algorithms:\\[4pt]
A Viability Assessment}}
\author{}
\date{May 2026}

\begin{document}
\maketitle
\vspace{-0.8em}

\section{The Proposal}

The system under consideration is DetailedBalanceChecker: a Mathematica symbolic
verifier that formally proves or disproves whether a Monte Carlo algorithm satisfies
detailed balance $T(i \to j)\,\pi(i) = T(j \to i)\,\pi(j)$.  It enumerates all
random-bit sequences via BFS, constructs the symbolic transition matrix, and tests
the balance condition algebraically---giving a formal YES/NO proof, not a
statistical test.

The proposal is to close a generative loop around this verifier: a large
language model (LLM) writes a candidate algorithm as Mathematica code; the
checker runs; the binary proof result and, on failure, the symbolic witness of
imbalance, are fed back to the LLM; the LLM revises and retries.  The goal is to
discover new Monte Carlo algorithms for lattice particle/polymer systems that
(a)~satisfy detailed balance by formal proof and (b)~exhibit dynamics
qualitatively closer to real polymer physics than existing lattice algorithms such
as Virtual Move Monte Carlo (VMMC) \cite{whitelam2007}.  This report assesses
where the proposal stands in the literature, identifies critical risks, and
describes what a working prototype would look like.

\section{Existing Algorithms and Correctness}

\textbf{Standard lattice algorithms.}  Kawasaki dynamics \cite{kawasaki1966}
(nearest-neighbour particle--hole swap with Metropolis acceptance) satisfies
detailed balance trivially: the proposal is symmetric and the Metropolis criterion
enforces balance.  VMMC \cite{whitelam2007} grows clusters by recruiting
neighbours with link probabilities
$w_\text{fwd} = \max(1 - e^{\beta(e_\text{init} - e_\text{fwd})}, 0)$; a
frustrated link aborts the move.  This enforces \emph{superdetailed balance}---a
condition strictly stronger than detailed balance---and the checker verifies it
analytically.  VMMC approximates overdamped Brownian dynamics heuristically;
its kinetics are physically plausible but not derived from the polymer equations
of motion.  Structurally related cluster algorithms include Swendsen--Wang
\cite{sw1987} and the Wolff algorithm \cite{wolff1989} for spin systems, and the
Liu--Luijten geometric cluster algorithm \cite{liu2004} for off-lattice fluids.

\textbf{Polymer Monte Carlo.}  The reptation family originates with de~Gennes'
tube theory \cite{degennes1971}.  The slithering-snake (Kron) move removes a
bond from one chain end and appends a new bond at the other end; it satisfies
detailed balance by construction (uniform bond-vector distribution, symmetric
forward and reverse moves) but is \emph{non-ergodic} for self-avoiding walks on
low-coordination lattices---a fact rigorously established by Madras \&
Sokal \cite{madras1988} for the pivot context and confirmed empirically for
slithering-snake dynamics by Wittmer et al.\ \cite{wittmer2003}.  This is the
sharpest caution for the user's project: the checker proves detailed balance but
does not assess ergodicity, which is a separate necessary condition.

Configurational-Bias MC (CBMC; Siepmann \& Frenkel \cite{siepmann1992}) grows a
chain segment bead-by-bead with biased trial positions, correcting the bias via
the ratio of Rosenbluth weights in the acceptance probability.  It is the
mathematically cleanest available template for satisfying detailed balance in
chain-growth moves and is ergodic by construction.  The pivot algorithm
\cite{madras1988} applies a random lattice symmetry to the distal chain
portion at a randomly chosen pivot; it satisfies detailed balance trivially and is
provably ergodic on hypercubic lattices, but is inapplicable at dense packing.
Houdayer's wormhole move \cite{houdayer2002} constructs a non-local
polymer-specific move from reptation steps that is provably detailed-balance
satisfying and faster than CBMC by $O(N)$.  End-bridging and double-bridging MC
\cite{pant1995,karayiannis2002} cut and reconnect chains at interior segments;
their acceptance criteria are non-trivial precisely because detailed balance must
be preserved, a task the user's checker could validate automatically.

It is also important to note that detailed balance is \emph{sufficient} but not
\emph{necessary} for the Boltzmann stationary distribution.  The weaker global
balance condition is sufficient \cite{manousiouthakis1999}, and algorithms that
exploit this---such as event-chain Monte Carlo \cite{ecmc2009}, which satisfies
only global balance via ``lifting'' variables and achieves rejection-free updates
with faster dynamics---would \emph{correctly fail} the user's detailed-balance
checker.  This is a feature: it defines a precise search objective.  It also means
a strictly larger class of physically valid algorithms lies outside the user's
current search space, and those algorithms (ECMC variants) may be more efficient.

\section{Prior Art: LLM-Driven Algorithm Discovery and Formal Verification}

Table~\ref{tab:priorart} summarises the most relevant published systems.

\begin{table}[h]
\small
\centering
\renewcommand{\arraystretch}{1.25}
\begin{tabular}{p{3.5cm} p{2.2cm} p{4.5cm} p{2.0cm}}
\toprule
\textbf{System} & \textbf{Year} & \textbf{Key feature} & \textbf{Physics MC?} \\
\midrule
FunSearch \cite{funsearch2023} & 2023 & LLM + evolutionary empirical evaluator & No \\
AlphaEvolve \cite{alphaevolve2025} & 2025 & Full codebase evolution, empirical test suites & No \\
AI Scientist \cite{aiscientist2024} & 2024 & LLM full research loop; empirical ML benchmarks & No \\
SGA \cite{sga2024} & 2024 & LLM hypotheses + differentiable simulation & No \\
POISE \cite{poise2025} & 2025 & LLM discovers RL algorithms; consistency checks & No \\
Clover \cite{clover2024} & 2024 & LLM + Dafny SMT formal verification & No \\
LLMLOOP \cite{llmloop2025} & 2025 & LLM + iterative compile/test feedback & No \\
Norm.\ flows for LFT \cite{albergo2021} & 2021 & Neural network MCMC proposals for lattice QFT & Partial \\
\bottomrule
\end{tabular}
\caption{Most relevant published systems. None applies LLM generation to physical MC algorithm design with a formal symbolic correctness oracle.}
\label{tab:priorart}
\end{table}

FunSearch \cite{funsearch2023} is the closest methodological ancestor: an LLM
proposes Python functions; an automated evaluator scores them; the best are evolved.
It discovered new cap-set solutions and improved bin-packing heuristics.
AlphaEvolve \cite{alphaevolve2025} (2025) extends this to evolving entire codebases
and found the first improvement to Strassen matrix multiplication in 56 years---but
relies on empirical performance scoring, not formal physical correctness.  Neither
has been applied to Monte Carlo algorithm design.

In the formal-verification literature, Clover \cite{clover2024} and LLMLOOP
\cite{llmloop2025} demonstrate that LLM+symbolic-checker loops are effective in
software engineering contexts (Lean/Dafny proofs, static analysis).  A salient
finding: ``compile-execute-revise loops are insufficient for scientific computing,
where a program may run successfully without producing correct results''
\cite{llmloop2025}.  This directly motivates the user's symbolic checker over a
runtime test suite.  However, none of these systems has been applied to physical
simulation algorithm design.

The combination---LLM proposer, formal detailed-balance checker as oracle, binary
proof fed back to guide revision, targeting polymer/soft-matter physics---does not
appear in the published literature.  The gap is real.

\section{Critical Assessment}

\textbf{What is genuinely strong.}  The symbolic checker provides a qualitatively
better feedback signal than empirical tests.  A symbolic counterexample---the
specific state pair $(i,j)$ and the non-zero residual expression---is far more
informative to the LLM than a failed KL divergence test.  The existing verified
examples (Kawasaki, VMMC and variants) constitute ready-made few-shot examples that
teach the LLM both the required code structure and the physical intuitions that have
previously worked.  The abstract-parameter capability of the checker (coupling and
field functions as free symbols) means new algorithms are verified for all possible
coupling constants simultaneously.  The methodological novelty over FunSearch is
real: replacing an empirical performance evaluator with a formal mathematical oracle
is a qualitative upgrade in the correctness guarantee.

\textbf{Ergodicity is not checked.}  This is the most serious risk.  An algorithm
that satisfies detailed balance but is non-ergodic (as the slithering-snake is) will
pass the checker with a clean PASS while producing incorrect equilibrium sampling.
The slithering-snake precedent is a direct warning: the checker would certify it.
Any prototype must address ergodicity separately, either by restricting to move
families that are ergodic by construction (CBMC-style full-chain regrowth) or by
implementing an ergodicity check as a BFS over the transition graph.

\textbf{The search space is vast and unstructured.}  Unconstrained LLM generation
from a blank prompt will overwhelmingly produce syntactically valid but physically
incoherent code.  The correct approach is constrained search: give the LLM a
skeleton algorithm---a physically motivated structure with a specific missing or
modifiable component---and ask it to fill in that component.  The VMMC codebase
itself documents that human experts introduce subtle errors (expanded neighbour
shell, double-test problem on $L=2$, \texttt{Piecewise} vs \texttt{Max}) that the
checker catches.  LLM generation will fail in analogous and new ways, requiring many
iterations.  This is acceptable if the loop is fast; the checker's runtime of
$\sim\!50$--$100$\,s for VMMC at \texttt{MaxBitString=1111111111} is
short enough to iterate hundreds of times in a single session.

\textbf{The physical target is underspecified.}  ``Closer to base reality of
polymer interaction'' requires operationalisation.  What observable?  Mean-squared
displacement scaling ($\langle r^2(t) \rangle \propto t^{1/2}$ for Rouse,
$\propto t^{1/4}$ for reptation), stress relaxation modulus, or chain renewal time
scaling?  Without a measurable kinetic target, neither the LLM nor the investigator
can determine whether a passing algorithm is an improvement over VMMC.  Detailed
balance is necessary for correct equilibrium sampling; it says nothing about
kinetics.

\textbf{Competing approaches are strong.}  CBMC already provides a principled
mathematical template for any chain-growth move.  AlphaEvolve could in principle be
applied to this problem if given a suitable evaluator, with far more computational
resources.  The user should be prepared to articulate what physical move CBMC cannot
express that the target algorithm would.  Event-chain MC \cite{ecmc2009} achieves
rejection-free updates with correct equilibrium statistics and dynamics often closer
to Brownian motion; it lies outside the detailed-balance search space but may be the
more powerful tool for the physical goals stated.

\textbf{LLMs can contribute but are not magic.}  The value of the LLM in this loop
is primarily as a rapid hypothesis generator: it can propose syntactically valid code
variants quickly and incorporate the symbolic counterexample into its revision.  It
cannot reason reliably about whether a proposed move is ergodic, whether its
kinetics will match any particular dynamical regime, or whether the acceptance
criterion correctly cancels a sampling bias without being given the relevant
mathematical formula.

\section{Publication Viability and Prototype}

The gap is real and the minimum viable publication consists of three things:
a new algorithm discovered via the LLM--checker loop that (a) passes the formal
symbolic checker, (b) is demonstrably ergodic, and (c) exhibits a kinetic
observable not reproduced by VMMC or Kawasaki on the same lattice model.  The
methodological contribution---the loop itself---is novel and worth documenting,
but it will not carry a paper alone without the algorithmic result.  Suitable
venues include \textit{J.\ Chem.\ Phys.}, \textit{J.\ Chem.\ Theory Comput.},
\textit{Phys.\ Rev.\ E}, or a methods letter at the AI--physics interface
(\textit{Mach.\ Learn.: Sci.\ Technol.}).

A working prototype has four components: \textbf{(1)}~A structured few-shot
prompt containing two passing examples (Kawasaki, VMMC) in full, a physical
description of the target move class, the most recent code attempt, and the
checker's symbolic counterexample if the attempt failed.  \textbf{(2)}~Checker
execution in symbolic mode (\texttt{check.wls}, \texttt{Mode=Symbolic},
\texttt{MaxBitString=11111} or higher).  \textbf{(3)}~Failure-analysis
extraction: parse the checker output to identify the specific $(i,j)$ pair and
residual expression, and include these verbatim in the next prompt.
\textbf{(4)}~Ergodicity verification: for each PASS result, run a BFS over the
transition graph from all reachable states to confirm irreducibility.

The most productive immediate target is a CBMC-style lattice chain-regrowth
move: delete a contiguous chain segment, regrow bead-by-bead with biased
sampling, and correct the bias via Rosenbluth weights in the acceptance
criterion.  The Rosenbluth weight formula can be given to the LLM explicitly as
a design pattern; the checker then verifies that the LLM's implementation
actually achieves cancellation.  This is ergodic by construction, physically
motivated (Siepmann \& Frenkel 1992 in the lattice context), and distinct from
VMMC in its kinetic character.

\section{An Extended Loop Architecture}

The preceding sections identified several gaps in the minimal
LLM--checker loop: ergodicity is not verified, there is no continuous
signal for kinetic quality, and the LLM is asked to guess acceptance
criteria that could instead be derived exactly.  This section develops
a concrete plan for a richer loop that addresses all three.

\subsection*{Ergodicity verification from the transition matrix}

The checker already constructs the full symbolic transition matrix
$T_{ij}$.  Ergodicity of the Markov chain is equivalent to strong
connectivity of the directed transition graph: every state must be
reachable from every other state.  The support matrix $S_{ij} = 1$ if
$T_{ij} \not\equiv 0$ (structurally nonzero, as established by
\texttt{=== 0} in Mathematica's symbolic algebra), and $S_{ij} = 0$
otherwise.  A standard result from graph theory guarantees that the
matrix power $(I + S)^{n-1}$ has a strictly positive $(i,j)$ entry if
and only if state $j$ is reachable from state $i$ in at most $n-1$
steps \cite{horn1985}.  The ergodicity check is therefore:

\begin{lstlisting}[language=Mathematica]
CheckErgodicity[T_List] := Module[{n, S, R},
  n = Length[T];
  (* Support: 1 if symbolically nonzero, 0 otherwise *)
  S = Map[If[# === 0, 0, 1] &, T, {2}];
  (* Reachability via matrix power (exact for integer matrices) *)
  R = MatrixPower[IdentityMatrix[n] + S, n - 1];
  If[Min[Flatten[R]] > 0,
    Print["ERGODIC: transition graph strongly connected."];
    True,
    (* Report the first unreachable pair *)
    With[{bad = First[Position[R, 0, {2}, 1]]},
      Print["NOT ERGODIC. State ", bad[[2]],
            " unreachable from state ", bad[[1]]]];
    False]
]
\end{lstlisting}

For the state-space sizes the checker examines (typically $n \lesssim
100$ states for small lattices), $n-1$ matrix multiplications of
integer matrices complete in milliseconds.  The check should be
inserted immediately after a PASS from the detailed-balance checker;
an algorithm that passes both is a valid candidate for further
evaluation.

The check gives \emph{generic} ergodicity: if $T_{ij}$ is structurally
nonzero but zero at a particular parameter value (e.g.\ $J = 0$),
the check correctly identifies the transition as generically available.
For degenerate parameter regimes, a separate numerical check at
problematic parameter values is advisable, but generic ergodicity is
the physically relevant condition for a non-trivial coupling constant.

\subsection*{Physical fidelity as a continuous optimisation signal}

Detailed balance certifies that the stationary distribution is
Boltzmann; ergodicity certifies that the chain reaches it.  Neither
condition says anything about whether the \emph{path} taken resembles
real physical dynamics.  A third signal is therefore needed: a scalar
score measuring how closely the MC transition matrix $T^{MC}$ matches
the physical transition matrix $T^{phys}$ derived from the actual
equations of motion of the modelled system.

\textbf{Computing $T^{phys}$ for the lattice model.}  For a lattice
gas with pairwise interactions, the physical dynamics is overdamped
Brownian motion projected onto lattice sites.  The transition rates
under this projection are given by Glauber dynamics \cite{glauber1963}:
for a single-particle hop from configuration $A$ to $B$ differing by
one lattice step,
\[
  k^{phys}_{A \to B} = \nu_0 \cdot
    \frac{e^{-\beta\Delta E/2}}{2\cosh(\beta\Delta E/2)},
\]
which is the symmetric Glauber rate (here $\Delta E = E_B - E_A$).
For cluster moves of $n$ particles under Stokes drag, the physical
rate scales as $k^{phys}_\text{cluster}(n) \propto D_n \propto 1/n$
in the appropriate dimension (with corrections for non-spherical
clusters).  For the small 2D lattice systems the checker examines,
$T^{phys}$ is a fully enumerable matrix: every pair of states connected
by a one-particle hop or a physically plausible multi-particle rearrangement
has a calculable rate.  Numerically this amounts to evaluating $\Delta E$
for each transition pair at specific parameter values $(\beta, J)$---a
negligible computation for the $n \lesssim 100$ state spaces involved.

\textbf{Comparing $T^{MC}$ and $T^{phys}$.}  A direct comparison of
matrix entries is confounded by the timescale ambiguity: the MC step
is not the same as the physical time unit.  The appropriate comparison
is therefore of the \emph{shapes} of the two matrices as row-conditional
distributions.  For each initial state $i$, define the normalised
physical rates $\tilde{T}^{phys}_{ij} = k^{phys}_{ij}/\sum_j
k^{phys}_{ij}$ and the MC proposal probabilities $\tilde{T}^{MC}_{ij}
= T^{MC}_{ij} / \sum_j T^{MC}_{ij}$.  The fidelity score is
\[
  \mathcal{F} = -\sum_i \pi_i \sum_j \tilde{T}^{phys}_{ij}
    \log \frac{\tilde{T}^{phys}_{ij}}{\tilde{T}^{MC}_{ij}},
\]
the negative KL divergence averaged over states weighted by the
equilibrium distribution.  Higher (less negative) is better; $\mathcal{F}
= 0$ means perfect shape match.  The score is undefined if $T^{MC}_{ij}
= 0$ for a transition with $\tilde{T}^{phys}_{ij} > 0$: the MC
algorithm cannot execute a physically possible transition.  This is a
hard failure and should be treated as disqualifying.

An alternative metric that avoids the undefined-log problem is
comparison of the non-trivial eigenvalue spectra:
$|\lambda^{MC}_k - c\,\lambda^{phys}_k|$ for a fitted scaling constant
$c$.  This measures whether the relaxation timescales of the two
dynamics have the same \emph{ratios}, even if their absolute scales
differ.  Matching the ratio $\lambda_2/\lambda_3$ (two slowest
non-trivial modes) is a stringent test of kinetic fidelity.

\textbf{Feasibility and practical implementation.}  For a 2$\times$2
lattice with 2 particle types, the full state space has $\sim$20--50
states, and all computations above complete in under a second.  For
3$\times$3 lattices ($\sim$1000 states), numerical evaluation of
$T^{phys}$ at a few parameter sets ($\beta \in \{0.5, 1.0, 2.0\}$,
$J \in \{-1, 0, +1\}$) takes minutes.  This is the right regime: the
physical fidelity signal should be evaluated numerically at a handful of
representative parameter values---not symbolically---and averaged.
The symbolic checker runs first (exact, parameter-free); the physical
fidelity scorer runs only on algorithms that pass both the DB check and
the ergodicity check.

\subsection*{The acceptance rate problem and rejection-free variants}

A specific and important inefficiency in standard VMMC is the following.
A cluster of $n$ particles can be initiated from any of its $n$ members,
so it is effectively attempted $n$ times per MC sweep.  Under Stokes drag
the physical diffusion coefficient of a compact cluster scales as $D_n
\propto 1/n$ in 3D (and analogously in 2D), so the physical displacement
rate is $\propto 1/n$.  For the MC attempt rate $(\propto n)$ times the
acceptance probability to equal the physical rate $(\propto 1/n)$, the
acceptance probability must scale as $\propto 1/n^2$.  This is a severe
penalty: for clusters of 10 particles the typical acceptance rate is
already $\sim$1\%, making large-cluster moves essentially decorative.

Note that this is distinct from the frustrated-link rejection, which is a
consequence of detailed balance and is in principle unavoidable given the
move structure.  The Stokes inefficiency is a consequence of the
\emph{proposal} distribution (uniform seed choice combined with fixed
step size) not matching the physical rate distribution.  It can in
principle be corrected without modifying the acceptance criterion.

\textbf{Proposal-level corrections.}  The natural fix is to \emph{weight
the seed selection inversely by cluster size}: assign seed particle $p$
a selection probability proportional to $1/n_p$ where $n_p$ is the size
of the cluster that would form from $p$.  This biases the algorithm
toward proposing small-cluster moves with the correct relative frequency.
The detailed-balance check must then be re-derived; the checker
would verify whether any specific implementation of this idea maintains
balance.  This is precisely the kind of modification the LLM is well
placed to propose: it is a small, local change to the seed-selection code
with a clear physical motivation, and the symbolic checker immediately
confirms or refutes whether balance is preserved.

\textbf{Automatically deriving the acceptance criterion.}  A more
radical approach exploits the symbolic checker directly.  Observe that
the checker already computes $q(i \to j)$---the probability of the
proposal kernel mapping state $i$ to state $j$---as a symbolic
expression.  Given $q$, the unique Metropolis acceptance probability
that satisfies detailed balance is
\[
  a(i \to j) = \min\!\left(1,\,
    \frac{q(j \to i)}{q(i \to j)}\,e^{-\beta\Delta E_{ij}}
  \right).
\]
Mathematica can compute the ratio $q(j \to i)/q(i \to j)$ symbolically
from the already-enumerated transition paths.  This means the workflow
can be restructured: the LLM proposes only the \emph{proposal
distribution} (the cluster growth rule, without an acceptance step);
Mathematica derives the acceptance criterion automatically and inserts
it; the checker then verifies the combined algorithm.  This separates
the creative task (proposing a physically motivated cluster growth rule)
from the mechanical task (computing the acceptance probability).  The
acceptance rate $\bar{A}(\beta, J)$ can then be computed symbolically
for small systems:
\[
  \bar{A}(\beta, J) = \sum_{i,j} \pi_i(\beta,J)\, q(i \to j)\,
    \min\!\left(1,\, \frac{q(j \to i)}{q(i \to j)}\,e^{-\beta\Delta E_{ij}}
    \right)
\]
and maximised over the free parameters in the growth rule.  This turns
algorithm discovery into an explicit optimisation problem with an exact
objective function, computable in Mathematica for any candidate proposal
the LLM generates.

The target is a growth rule for which $\bar{A}$ is high (approaching 1)
and $\mathcal{F}$ (the physical fidelity score defined above) is also
high.  These two objectives can conflict: a proposal that always accepts
may do so by making moves that are physically unmotivated, giving low
$\mathcal{F}$.  Managing this tension is where physical intuition (and
the LLM's knowledge of Brownian dynamics) becomes essential.

\subsection*{How an LLM should exploit Mathematica's symbolic engine}

The standard LLM--code loop (write code, run it, observe output, fix
code) uses the execution environment as an oracle for binary
correctness.  The present system can go further, because Mathematica is
a computer algebra system, not merely a runtime.  The LLM can interact
with it in several qualitatively richer ways.

\textbf{Receiving algebraic error witnesses.}  When the detailed-balance
checker fails, it does not merely say ``wrong''; it returns the specific
symbolic expression $T_{ij}\pi_i - T_{ji}\pi_j$ that is nonzero for
a specific state pair.  This expression is a polynomial in $e^{\beta
J_{ab}}$ terms.  The LLM can parse this expression and identify the
structural cause of the imbalance.  For example, if the residual is
$e^{\beta J_{12}} - e^{\beta J_{13}}$, the algorithm is treating bonds
of type (1,2) and (1,3) asymmetrically when it should treat them
identically.  This is a diagnostic far richer than any runtime error
message.

\textbf{Proposing and testing parametric families.}  Rather than
proposing a single algorithm, the LLM can propose a \emph{parametric
family} with free parameters (e.g.\ the frustration threshold, a
weighting exponent, a cutoff distance).  Mathematica computes the
acceptance rate and fidelity score as symbolic functions of these
parameters, and the optimal values can be found by \texttt{Maximize} or
\texttt{FindMaximum}.  This allows the LLM to propose general
algorithmic structures and let Mathematica tune them---a more productive
division of labour than asking the LLM to guess optimal parameter values.

\textbf{Symbolic consistency checks before running the full checker.}
For a candidate algorithm, certain necessary conditions for detailed
balance can be checked symbolically in microseconds, without running the
full BFS enumeration.  For example: are the transition probabilities
normalised (each row sums to 1)?  Is $T_{ij} = 0$ whenever $E_j -
E_i > E_\text{max}$ (energy monotonicity)?  Is $T_{ij}/T_{ji} =
e^{-\beta\Delta E_{ij}}$ for all directly connected pairs (local balance)?
The LLM can be instructed to include these checks in its generated code
as self-diagnostic assertions, catching structural errors before the
expensive BFS enumeration runs.

\textbf{Guided mutation of the acceptance criterion.}  Once the symbolic
ratio $q(j \to i)/q(i \to j)$ is computed, the LLM can be prompted:
``the current acceptance ratio is $R = [\text{expression}]$; propose a
modification to the growth rule that makes $R$ closer to 1 for large
clusters while keeping the expression valid as a probability''.  This
is a well-posed algebraic question that the LLM can reason about in
symbolic terms, guided by the exact expression from Mathematica rather
than a vague description of the imbalance.

\subsection*{Connection to embedded LLM theory}

The theoretical framework developed in the companion document
\texttt{embedded\_LLM.tex} \cite{gao2023pal,yao2023react,shinn2023reflexion}
identifies three mechanisms by which embedding an LLM in an executable
environment improves performance: unambiguous external feedback (the
environment cannot be fooled by fluent text), offloading of formal
computation (the interpreter does exactly what the LLM approximates
poorly), and error localisation (failures are pinpointed rather than
compounding).

The present system is a particularly rich instantiation of all three.
The detailed-balance checker provides feedback the LLM cannot
hallucinate: the algebraic expression $T_{ij}\pi_i - T_{ji}\pi_j$ is a
theorem, not a probability.  The acceptance rate $\bar{A}(\beta,J)$ and
fidelity score $\mathcal{F}$ are formal computations that Mathematica
evaluates exactly in seconds and the LLM could not reliably compute
internally.  And the symbolic counterexample localises the error to a
specific state pair and a specific algebraic term.

The proposed system adds a \emph{fourth} mechanism not discussed in the
embedded-LLM literature: \textbf{symbolic reasoning feedback}.  A Python
traceback tells the LLM ``your code crashed at line 12''.  The symbolic
checker tells the LLM ``the detailed balance condition fails for states
$(3, 7)$ because the expression $e^{\beta J_{12}} - e^{\beta J_{13}}$ is
nonzero''.  This is not a runtime error; it is a mathematical statement
about the structure of the algorithm.  The LLM can reason about it
algebraically---recognising which part of the growth rule generated the
$J_{12}$ term and which generated the $J_{13}$ term, and modifying the
code to symmetrise them.  This is qualitatively beyond what execution
feedback in standard programming contexts can provide: the oracle speaks
in the same algebraic language as the problem domain.

This feature also resolves a key limitation identified in the embedded-LLM
literature: that ``execution loops are insufficient for scientific
computing, where a program may run correctly without producing correct
results'' \cite{llmloop2025}.  The symbolic checker does not test
outcomes; it tests the underlying mathematical relationship.  A program
that runs without error and produces numerically plausible-looking output
but violates detailed balance will be identified as incorrect---exactly
the failure mode that runtime tests cannot catch.  The combination of
three independent formal verifiers (detailed balance, ergodicity,
physical fidelity) provides a layered guarantee that no existing LLM
algorithm discovery system approaches.

\section{Conclusion}

The proposal is scientifically sound and fills a genuine gap.  The core
innovation---using a formal mathematical oracle in place of an empirical evaluator
in an LLM discovery loop---is both conceptually sharp and practically feasible with
the existing checker infrastructure.  The principal risks are the unchecked
ergodicity problem, the underspecification of the physical target, and the breadth
of the search space.  None is fatal; all require explicit mitigation.  If the LLM
loop produces one new algorithm that passes the formal checker, is ergodic, and
distinguishes itself kinetically from VMMC, there is a credible publication.

\begin{thebibliography}{99}

\bibitem{whitelam2007}
  Whitelam, S.\ \& Geissler, P.\,L.\ (2007).
  Avoiding unphysical kinetic traps in Monte Carlo simulations of strongly
  attractive particles.
  \textit{J.\ Chem.\ Phys.}\ \textbf{127}, 154101.

\bibitem{kawasaki1966}
  Kawasaki, K.\ (1966).
  Diffusion constants near the critical point for time-dependent Ising models.
  \textit{Phys.\ Rev.}\ \textbf{145}, 224.

\bibitem{sw1987}
  Swendsen, R.\,H.\ \& Wang, J.-S.\ (1987).
  Nonuniversal critical dynamics in Monte Carlo simulations.
  \textit{Phys.\ Rev.\ Lett.}\ \textbf{58}, 86.

\bibitem{wolff1989}
  Wolff, U.\ (1989).
  Collective Monte Carlo updating for spin systems.
  \textit{Phys.\ Rev.\ Lett.}\ \textbf{62}, 361.

\bibitem{liu2004}
  Liu, J.\ \& Luijten, E.\ (2004).
  Rejection-free geometric cluster algorithm for complex fluids.
  \textit{Phys.\ Rev.\ Lett.}\ \textbf{92}, 035504.

\bibitem{degennes1971}
  de~Gennes, P.\,G.\ (1971).
  Reptation of a polymer chain in the presence of fixed obstacles.
  \textit{J.\ Chem.\ Phys.}\ \textbf{55}, 572.

\bibitem{madras1988}
  Madras, N.\ \& Sokal, A.\,D.\ (1988).
  The pivot algorithm: a highly efficient Monte Carlo method for the
  self-avoiding walk.
  \textit{J.\ Stat.\ Phys.}\ \textbf{50}, 109.

\bibitem{siepmann1992}
  Siepmann, J.\,I.\ \& Frenkel, D.\ (1992).
  Configurational bias Monte Carlo: a new sampling scheme for flexible chains.
  \textit{Mol.\ Phys.}\ \textbf{75}, 59.

\bibitem{houdayer2002}
  Houdayer, J.\ (2002).
  The wormhole move: a new algorithm for polymer simulations.
  \textit{J.\ Chem.\ Phys.}\ \textbf{116}.
  arXiv:cond-mat/0111087.

\bibitem{pant1995}
  Pant, P.\,V.\,K.\ \& Theodorou, D.\,N.\ (1995).
  Variable connectivity method for the atomistic Monte Carlo simulation of
  polydisperse polymer melts.
  \textit{Macromolecules}\ \textbf{28}, 7224.

\bibitem{karayiannis2002}
  Karayiannis, N.\,C., Mavrantzas, V.\,G.\ \& Theodorou, D.\,N.\ (2002).
  A novel Monte Carlo scheme for the rapid equilibration of atomistic model
  polymer systems.
  \textit{Phys.\ Rev.\ Lett.}\ \textbf{88}, 105503.

\bibitem{wittmer2003}
  Wittmer, J.\,P.\ et al.\ (2003).
  Dynamical properties of the slithering-snake algorithm: a numerical test
  of the activated-reptation hypothesis.
  \textit{Eur.\ Phys.\ J.\ E}\ \textbf{10}, 369.
  arXiv:cond-mat/0212433.

\bibitem{manousiouthakis1999}
  Manousiouthakis, V.\,I.\ \& Deem, M.\,W.\ (1999).
  Strict detailed balance is unnecessary in Monte Carlo simulation.
  \textit{Phys.\ Rev.\ E}\ \textbf{59}, 2.

\bibitem{ecmc2009}
  Bernard, E.\,P., Krauth, W.\ \& Wilson, D.\,B.\ (2009).
  Event-chain Monte Carlo algorithms for hard-sphere systems.
  \textit{Phys.\ Rev.\ E}\ \textbf{80}, 056704.

\bibitem{funsearch2023}
  Romera-Paredes, B.\ et al.\ (2023).
  Mathematical discoveries from program search with large language models.
  \textit{Nature}\ \textbf{628}, 73.

\bibitem{alphaevolve2025}
  Google DeepMind (2025).
  AlphaEvolve: a coding agent for scientific and algorithmic discovery.
  arXiv:2506.13131.

\bibitem{aiscientist2024}
  Lu, C.\ et al.\ (2024).
  The AI Scientist: towards fully automated open-ended scientific discovery.
  arXiv:2408.06292.

\bibitem{sga2024}
  Ma, P.\ et al.\ (2024).
  LLM and simulation as bilevel optimizers: a new paradigm to advance
  physical scientific discovery.
  \textit{ICML 2024}. arXiv:2405.09783.

\bibitem{poise2025}
  From AI assistant to AI scientist: autonomous discovery of LLM-RL
  algorithms with LLM agents. (2025). arXiv:2603.23951.

\bibitem{clover2024}
  Clover: closed-loop verifiable code generation. (2024, AAAI adjacent).

\bibitem{llmloop2025}
  LLMLOOP: automating refinement of LLM-generated code through iterative
  feedback. \textit{ICSME} (2025). arXiv:2603.23613.

\bibitem{albergo2021}
  Albergo, M.\,S.\ et al.\ (2021).
  Introduction to normalizing flows for lattice field theory.
  arXiv:2101.08176.

\bibitem{horn1985}
  Horn, R.\,A.\ \& Johnson, C.\,R.\ (1985).
  \textit{Matrix Analysis}.  Cambridge University Press.

\bibitem{glauber1963}
  Glauber, R.\,J.\ (1963).
  Time-dependent statistics of the Ising model.
  \textit{J.\ Math.\ Phys.}\ \textbf{4}, 294.

\bibitem{gao2023pal}
  Gao, L.\ et al.\ (2023).
  PAL: Program-aided language models.
  \textit{Proc.\ ICML}, PMLR 202:10764.

\bibitem{yao2023react}
  Yao, S.\ et al.\ (2023).
  ReAct: Synergizing reasoning and acting in language models.
  \textit{ICLR 2023}.

\bibitem{shinn2023reflexion}
  Shinn, N.\ et al.\ (2023).
  Reflexion: Language agents with verbal reinforcement learning.
  \textit{NeurIPS 2023}.

\end{thebibliography}

\end{document}
